\documentclass[10pt,a4paper]{amsart}
\usepackage[margin=19mm]{geometry}
\usepackage{amsmath,amssymb,mathtools}
\usepackage[scr=rsfs,cal=euler]{mathalfa}
\usepackage{xfrac}
\usepackage[unicode,hidelinks,bookmarks=false]{hyperref}
\newcommand{\ie}{i.e.\ }
\newcommand{\cf}{cf.\ }
\newcommand{\resp}{respectively\ }
\newcommand{\Subsection}{\subsection}
\providecommand{\nobreakdash}{\nobreak\hbox{-}}
\setlength{\emergencystretch}{2em}
\multlinegap0pt
\usepackage{bm}
\usepackage{bbm}
\usepackage{dsfont}
\usepackage{xspace}
\usepackage{cancel}
\usepackage{mathtools}
\usepackage{amsthm}
\makeatletter
\@ifundefined{thm}{\newtheorem{thm}{Theorem}}{}
\@ifundefined{lem}{\newtheorem{lem}[thm]{Lemma}}{}
\makeatother
\def\diam{{\mathop\mathrm{\,diam\,}}}
\def\dist{{\mathop\mathrm{\,dist\,}}}
\def\ls{\lesssim}
\title[Planar $W^{1,\,1}$-extension domains]{Planar $W^{1,\,1}$-extension domains}
\author{Pekka Koskela, Tapio Rajala, Yi Ru-Ya Zhang}
\date{Source: arXiv:2512.09167v2 (CC BY 4.0)}
\begin{document}
\maketitle

\noindent\emph{Source and licence.} Planar $W^{1,\,1}$-extension domains. Version arXiv:2512.09167v2 (CC BY 4.0), distributed under the Creative Commons Attribution 4.0 International licence, as recorded in the arXiv licence metadata for this exact version and shown on the abstract page https://arxiv.org/abs/2512.09167v2. Full rights notice: licenses/arxiv-2512.09167-CC-BY-4.0.txt.\par\smallskip

\noindent\textit{Editorial scope.} Counted unit: the Lem whose source label is \texttt{inner capacity} in the article (source0.tex lines 2908--2921, source label \texttt{inner capacity}), delivered together with the article's own complete original proof of it (source0.tex lines 2922--3048, one \texttt{proof} environment of 127 lines). No mathematics is written, repaired or re-derived: statement and proof are the article's own text line for line. Required context retained uncounted: GM bdd (lines 500--508). Every definition, display and label of those lines is retained unchanged. Omitted from this package and not redistributed: 1--499, 509--2907, 3049--3429. Nothing inside a retained excerpt is omitted or altered: every retained body is delivered exactly as the article's lines are written, so no omission receipt of any kind accompanies this package. Citations: the retained excerpts carry no citation command at all, so no bibliography entry of the article is transcribed and no reference is redistributed. Reference adaptation: none was needed and none was made; every reference inside the delivered unit resolves inside it, so no unresolved reference or citation is produced. Printed slips kept literal: no printed word, symbol, hypothesis or number of the counted statement or of its proof was altered. Layout and class: the article's own class is \texttt{amsart} with packages including a4wide, color, eucal, enumerate, mathrsfs, ulem, amsmath, amssymb, epsfig, amsthm, inputenc, psfrag, hyperref, cleveref; the wrapper is the lane's pinned \texttt{amsart} editorial wrapper at 10pt on A4, which restates the theorem-like environments the retained text needs and adds the article's own macro definitions that the retained text needs (\texttt{\textbackslash diam}, \texttt{\textbackslash dist}, \texttt{\textbackslash ls}), with extra wrapper packages (bm, bbm, dsfont, xspace, cancel, mathtools). The delivered numbering is the unit's own. Assets: no retained span contains a figure environment, a graphics-inclusion command, a PDF-page inclusion, a table or a TikZ picture; no image file is copied into this package, the package declares no asset dependency and no image file is redistributed with it. Licence: version arXiv:2512.09167v2 of this article is distributed under the Creative Commons Attribution 4.0 International licence, as recorded in the arXiv licence metadata for this exact version and shown on the abstract page https://arxiv.org/abs/2512.09167v2. A full rights notice is delivered at licenses/arxiv-2512.09167-CC-BY-4.0.txt. Characters: every retained excerpt is pure ASCII, so the delivered engine, XeLaTeX with its default Latin Modern text font, reports no missing character.
\begin{thm}\label{GM bdd}
Let $\Omega$ be a bounded simply connected domain, and $\varphi\colon \mathbb D \to \Omega$ be a conformal map. Then there exists an absolute constant $C>0$ so that,  given a pair of points
$x,\,y\in \mathbb D$, 
denoting the corresponding hyperbolic geodesic in 
$\mathbb D$ by $\Gamma_{x,\,y}$, 
and by $\gamma_{x,\,y}$ any 
arc connecting $x$ and $y$ in $ \mathbb D$, we have
$$\ell(\varphi(\Gamma_{x,\,y}))\le C \ell(\varphi(\gamma_{x,\,y})). $$
\end{thm}

\begin{lem}{\label{inner capacity}}
Let $\Omega$ be a domain and  $E,\,F\subset \Omega$ be a pair of disjoint continua. 
Then if ${\rm Cap}(E,\,F,\, \Omega)\ge c_0$, we have
\begin{equation}\label{inequat 1}
{\min\{\diam_{\Omega}(E),\,\diam_{\Omega}(F)\}}\gtrsim {\dist_{\Omega}(E,\,F)},
\end{equation}
where the constant only depends on $c_0$. Especially
$${\min\{\diam_{\Omega}(E),\,\diam_{\Omega}(F)\}}\gtrsim {\dist (E,\,F)}, $$
and if $\Omega=\mathbb R^2$
\begin{equation}\label{inequat 10}
\min\{\diam(E),\,\diam(F)\}\gtrsim {\dist (E,\,F)}.  
\end{equation}
If we further assume that $\Omega$ is Jordan, then \eqref{inequat 1} also holds if $E\subset \partial\Omega$ and $F\subset\Omega$ (or $E\subset \partial\Omega$ and $F\subset\partial\Omega$) are continua with ${\rm Cap}(E,\,F,\, \Omega)\ge c_0$. 
\end{lem}

\begin{proof}
\noindent{\bf Step 1:}
We begin with the case where $E,\,F\subset \Omega$. 
We may clearly assume that $\diam_{\Omega}(E)\le \diam_{\Omega}(F)$ and also that
$2\diam_{\Omega}(E) \le \dist_{\Omega}(E,\,F)$; otherwise the claim holds trivially. 
Choose $z\in E$, and write $\frac{\dist_{\Omega}(E,\,F)}{\diam_{\Omega}(E)}=\sigma$. 
We define
$$ u(x)=
 \begin{cases}
   1, & \text{if } \dist_{\Omega}(x,\,z)\le \diam_{\Omega}(E) \\
  0, & \text{if}  \dist_{\Omega}(x,\,z)\ge \dist_{\Omega}(E,\,F)\\
      \frac {\log (\dist_{\Omega}(E,\,F))-\log (\dist_{\Omega}(x,\,z))}{\log (\dist_{\Omega}(E,\,F))-\log (\diam_{\Omega}(E))}, &\text{otherwise }
 \end{cases}.$$

Then $u$ is locally Lipschitz and 
$$|\nabla u(x)|\le (\log \sigma)^{-1} \dist_{\Omega}(x,\,z)^{-1}.$$
Write
$$R=B_{\Omega}(z,\,\dist_{\Omega}(E,\,F))\setminus B_{\Omega}(z,\,\diam_{\Omega}(E)),$$
and for $i\ge 1$
$$A_i=B_{\Omega}(z,\,2^{i}\diam(E))\setminus B_{\Omega}(z,\,2^{i-1}\diam(E)),$$
where $B_{\Omega}(z,\,r)$ is the disk centered at $z$ with radius $r$ with respect to the inner distance. 
A direct calculation via a dyadic annular decomposition with respect
to the inner distance gives
\begin{align*}
c_0\le & \int_{\Omega} |\nabla u|^2\, dx  
\ls   (\log \sigma)^{-2}\int_{R} \dist_{\Omega}(x,\,z)^{-2} \, dx\\ 
\ls &   (\log \sigma)^{-2} \sum_{i=1}^{\infty}\int_{R\cap A_i } 4^{-i}\diam(E)^{-2} \, dx\\
\ls &   (\log \sigma)^{-2} \sum_{i=1}^{[\log \sigma]+1} 1\\
\ls &  (\log \sigma)^{-2} \log \sigma \lesssim  {(\log \sigma)}^{-1}, 
\end{align*}
where  $[\log \sigma]$ denotes the integer part of $[\log \sigma]$, and in the fourth inequality we used the fact that $B_{\Omega}(z,\,r)\subset B(z,\,r)$. 
Hence $\sigma\le C(c_0)$, which means that $\dist_{\Omega}(E,\,F) \lesssim \diam_{\Omega}(E)$. 


\noindent{\bf Step 2:} We continue with the case $E,\,F\subset\partial\Omega$. 
To begin with, let us first show that 
$$\dist_{\Omega}(E,\,F)<\infty. $$
Let $\varphi\colon \overline{\mathbb D} \to \overline{\Omega}$ be a homeomorphism given by the Caratheodory-Osgood theorem. 
Since 
$${\rm Cap}(E,\,F,\, \mathbb R^2)\ge {\rm Cap}(E,\,F,\, \Omega)\ge c_0,$$
we conclude by \eqref{inequat 10} that neither $E$ nor $F$ is a singleton. 
Then $\varphi^{-1}(E)$ is an closed arc contained in the unit circle and of positive $1$-Hausdorff measure. Since $\varphi\in W^{1,\,2}(\mathbb D)$, then
$$\dist_{\Omega}(E,\,\varphi(0))<\infty;$$
otherwise by noticing that $\varphi$ is a smooth map on $\overline{B(0,\,1/2)}$ and applying Fubini's theorem  in the annulus $\mathbb D \setminus \overline{B(0,\,1/2)}$ with respect to the polar coordinates, we obtain a contradiction to the fact that $\varphi\in W^{1,\,2}(\mathbb D)$. 
An analogous argument also shows that 
$$\dist_{\Omega}(F,\,\varphi(0))<\infty,$$
and hence by the triangle inequality $\dist_{\Omega}(E,\,F)<\infty. $



Let $\phi\colon \mathbb H\to \Omega$ be a conformal map (homeomorphically extended  to  $\mathbb R$), where $\mathbb H$ denotes the upper half plane so that $\phi^{-1}(E)$ and $\phi^{-1}(F)$ are two continua in the real line. 
We define a sequence of continua $E_j \subset  \phi^{-1}(E)$ (for large $j$) by setting
$$ E_j =\{ z \in \mathbb R \colon   \dist(z,\,\mathbb R \setminus \phi^{-1}(E))\ge 2^{-j+1}\}. $$
The sets $ F_j\subset \phi^{-1}(F)$ are defined similarly. 
Furthermore, define
$$\hat E_j= E_j +i  2^{-j}, \ \hat F_j= F_j +i 2^{-j}. $$


Note that for every two points $z_1,\,z_2\in E$, the hyperbolic geodesic $\Gamma'$ joining them satisfies (see Lemma~\ref{GM bdd})
$$\ell(\Gamma')\lesssim \dist_{\Omega}(z_1,\,z_2), $$
where the constant is absolute. 
Given $w_1,\,w_2\in \phi(\hat E_j)$, denote the hyperbolic geodesic connecting them by $\Gamma$. Extend $\Gamma$ to a full hyperbolic geodesic $\Gamma_1$ joining points $z_1,\,z_2\in \partial \Omega$. Since $w_1,\,w_2\in \phi(\hat E_j)$, from the definition of $\hat E_j$ and planar geometry  we conclude via $\phi$ that $z_1,\,z_2\in  E$. Therefore
$$\dist_{\Omega}(w_1,\,w_2)\le \ell(\Gamma)\le \ell(\Gamma_1)\lesssim \dist_{\Omega}(z_1,\,z_2)\le \diam_{\Omega}(E). $$
By the arbitrariness of $w_1,\,w_2$, we conclude that
\begin{equation}\label{inequ 19}
\diam_{\Omega}(\phi(\hat E_j)) \lesssim  \diam_{\Omega}(E) 
\end{equation}
with an absolute constant. Similarly we get  
\begin{equation}\label{inequ 20}
\diam_{\Omega}(\phi(\hat F_j)) \lesssim  \diam_{\Omega}(F). 
\end{equation}

By the monotonicity of capacity, we have
$${\rm Cap}(\hat E_j,\,\hat F_j,\, \mathbb H)\ge {\rm Cap}(\hat E_j,\,\hat F_j,\, \mathbb H + i 2^{-j}  ). $$ 
Via the translation map $z\mapsto  z - i2^{-j}$ and the conformal invariance of capacity, we further have
$${\rm Cap}(\hat E_j,\, \hat F_j,\, \mathbb H + i 2^{-j} )={\rm Cap}(E_j,\, F_j,\,  \mathbb H). $$
 If we can show that
\begin{equation}\label{inequ 21}
{\rm Cap}(E_j,\, F_j,\,  \mathbb H)\ge \frac 1 {16} {\rm Cap}(\phi^{-1}(E) ,\, \phi^{-1}(F) ,\,  \mathbb H)
\end{equation}
for $j$ large enough, then by the conformal invariance of capacity together with the chain of inequalities above
$$ {\rm Cap}(\phi(\hat E_j),\,\phi(\hat F_j),\,\Omega )\ge \frac 1 {16} {\rm Cap}(\phi^{-1}(E) ,\, \phi^{-1}(F) ,\,  \mathbb H)= \frac 1 {16} {\rm Cap}( E ,\,  F  ,\,  \Omega)\ge \frac {c_0} {16}.$$
Thus Lemma~\ref{inner capacity} implies that
$${\min\{\diam_{\Omega}(\phi(\hat E_j)),\,\diam_{\Omega}(\phi(\hat F_j))\}}\gtrsim {\dist_{\Omega}(\phi(\hat E_j),\,\phi(\hat F_j))}, $$
where the constant only depends on $c_0$.
This together with \eqref{inequ 19} and  \eqref{inequ 20} gives that there is a rectifiable curve 
$\gamma_j$ connecting $\phi(\hat E_j)$ and $\phi(\hat F_j)$ such that
$$\ell(\gamma_j)\lesssim {\min\{\diam_{\Omega}(\phi(\hat E_j)),\,\diam_{\Omega}(\phi(\hat F_j))\}}\ls \min\{\diam_{\Omega}(E),\,\diam_{\Omega}(F)\},$$
with the constant independent of $j$. 
We may assume that these curves are hyperbolic geodesics by Lemma~\ref{GM bdd}.
Then by parameterizing $\gamma_j$ with arc length and applying the Arzel\'a-Ascoli lemma, we get a  curve $\gamma\subset \overline{\Omega}$ joining $E,\,F$ with the length bounded by $\min\{\diam_{\Omega}(E),\,\diam_{\Omega}(F)\}$ up to a multiplicative constant. Observe that, since $\varphi$ is a homeomorphism, the uniform convergence of $\gamma_j\to \gamma$ implies that of $\varphi^{-1}(\gamma_j)\to \varphi^{-1}(\gamma)$. Then the definition of hyperbolic geodesics in the unit disk gives that $\varphi^{-1}(\gamma)$ is also a hyperbolic geodesic, and so is  $\gamma$. 
The desired estimate \eqref{inequat 1} then follows. 
Thus it suffices to show \eqref{inequ 21}. 



 Set 
$$a=\min\{(\diam(\phi^{-1}(E)),\,\diam(\phi^{-1}(F)),\,\dist(\phi^{-1}(E),\,\phi^{-1}(F))\}.$$ 
Let the middle point of $\phi^{-1}(E)$ be $z_1$ and that of $\phi^{-1}(F)$ be $z_2$. 
Write 
\begin{equation}\label{inequat 2}
b_E=\frac 1 2 \diam(\phi^{-1}(E)) +\frac a 8,\ b_F=\frac 1 2 \diam(\phi^{-1}(F)) +\frac a 8.
\end{equation}
Recall that
$$\diam(\phi^{-1}(E))-\diam(E_j) = 2^{-j+2},\ \diam(\phi^{-1}(F))-\diam(F_j) = 2^{-j+2}. $$
For $j\in \mathbb N,\, 2^{-j+5}\le a$, we define a map $f$ by setting
$$ f(x)=
 \begin{cases}
   z_1+\frac{\diam(\phi^{-1}(E))}{\diam(E_j)}(x-z_1), & \,  |x-z_1|\le \frac 1 2 \diam(E_j) \\
  z_1+\left(\frac 1 2 \diam(\phi^{-1}(E)) + \frac{ a \left(|x-z_1|-\frac 1 2 \diam(E_j)\right)}{ 2^{-j+4}+ a}  \right)\frac {(x-z_1)}{|x-z_1|}, &   \, \frac 1 2 \diam(E_j)\le |x-z_1|\le b_E \\
   z_2+\frac{\diam(\phi^{-1}(F))}{\diam(F_j)}(x-z_2), & \,  |x-z_2|\le \frac 1 2 \diam(F_j) \\
   z_2+\left(\frac 1 2 \diam(\phi^{-1}(F)) + \frac{ a \left(|x-z_2|-\frac 1 2 \diam(F_j)\right) }{ 2^{-j+4}+ a}\right) \frac {(x-z_2)}{|x-z_2|}, &   \, \frac 1 2 \diam(F_j)\le |x-z_2|\le b_F \\
    x , & \text{ otherwise.}
 \end{cases}$$
 Then $f$ is a well-defined homeomorphism; recall \eqref{inequat 2} and notice that
 $$b_E-\frac 1 2 \diam(E_j)= 2^{-j+1}+\frac a 8, \ \ b_F-\frac 1 2 \diam(F_j)= 2^{-j+1}+\frac a 8.$$
By our choice of $j$, $f$  maps $\mathbb H$ to $\mathbb H$, $E_j$ to $\phi^{-1}(E)$ and $F_j$ to $\phi^{-1}(F)$. Moreover since $2^{-j+5}\le a$, 
$$ \frac 1 2 \le \frac { a }{ 2^{-j+4}+a}\le 1 ,\ 1\le \frac{\diam(\phi^{-1}(E))}{\diam(E_j)}\le 2,\  1\le \frac{\diam(\phi^{-1}(F))}{\diam(F_j)} \le 2.$$
Thus $f$ is $2$-biLipschitz. For any $u\in \Delta(E_j ,\, F_j,\,\mathbb H)$ we have $u\circ f^{-1}\in \Delta( \phi^{-1}(E) ,\,  \phi^{-1}(F),\,\mathbb H)$ since $f$ is a homeomorphism. Thus by the chain rule and  a change of variables we get
\begin{align*}
{\rm Cap}(\phi^{-1}(E) ,\, \phi^{-1}(F) ,\,  \mathbb H)& \le  \int_{\mathbb H} |\nabla u(f^{-1}(x))|^2 |Df^{-1}(x)|^2 \, dx\\
&\le 4 \int_{\mathbb H} |\nabla u(x)|^2  J_{f}(x) \, dx \le 16 \int_{\mathbb H} |\nabla u|^2 \, dx,
\end{align*}
which implies \eqref{inequ 21} as desired, and then \eqref{inequat 1} follows. 

\noindent{\bf Step 3:} We are left with the case  where $E\subset \partial\Omega$ and $F\subset\Omega$. In this case we  only perform the above approximation for $E$ (and consider $j$ large enough so that $E_j\cap \phi^{-1}(F)=\emptyset$). The argument in the previous case applies with this modification and gives the remaining claim. 
\end{proof}

\end{document}
